%%%%%%%%%%%%%%%%%%%%%%%%%%%%%%%%%%%%%%%%%%%%%%%%%%%%%%%%%%
\begin{frame}
  \begin{center}
    {\Large End-to-End Example: Customer Churn}
  \end{center}
\end{frame}

%%%%%%%%%%%%%%%%%%%%%%%%%%%%%%%%%%%%%%%%%%%%%%%%%%%%%%%%%%
\begin{frame}[fragile]\frametitle{The Question}
A telecom company wants to know: which customers are likely to \emph{churn} (cancel their
service)? We'll walk one dataset through the full arc -- load, assess, describe, visualize,
engineer features, check what matters -- the same pattern you'll reuse on almost every
dataset in this course.
\begin{itemize}
\item Dataset: per-customer usage records (call minutes, charges, plan type) plus whether
  that customer eventually churned.
\end{itemize}
\end{frame}

%%%%%%%%%%%%%%%%%%%%%%%%%%%%%%%%%%%%%%%%%%%%%%%%%%%%%%%%%%
\begin{frame}[fragile]\frametitle{Step 1: Load Data}
\begin{lstlisting}
import pandas as pd

df = pd.read_csv('telecom_churn.csv')
df.shape
(3333, 20)
\end{lstlisting}
\end{frame}

%%%%%%%%%%%%%%%%%%%%%%%%%%%%%%%%%%%%%%%%%%%%%%%%%%%%%%%%%%
\begin{frame}[fragile]\frametitle{Step 2: Initial Assessment}
Before anything else: what columns exist, and what type is each one?
\begin{lstlisting}
df.columns
Index(['State', 'Account length', 'Area code', 'International plan',
       'Voice mail plan', 'Total day minutes', 'Total day calls',
       'Total eve minutes', 'Customer service calls', 'Churn', ...])

df.dtypes
State                object
International plan  object
Total day minutes    float64
Customer service calls  int64
Churn                  bool
dtype: object
\end{lstlisting}
\end{frame}

%%%%%%%%%%%%%%%%%%%%%%%%%%%%%%%%%%%%%%%%%%%%%%%%%%%%%%%%%%
\begin{frame}[fragile]\frametitle{Step 3: Descriptive Statistics}
\texttt{describe()} gives a fast univariate read on every numeric column at once -- range,
center, spread.
\begin{lstlisting}
df.describe()
       Account length  Total day minutes  Customer service calls
count         3333.00            3333.00                  3333.00
mean           101.06             179.78                     1.56
std             39.82              54.47                     1.32
min              1.00               0.00                     0.00
max            243.00             350.80                     9.00
\end{lstlisting}
\end{frame}

%%%%%%%%%%%%%%%%%%%%%%%%%%%%%%%%%%%%%%%%%%%%%%%%%%%%%%%%%%
\begin{frame}[fragile]\frametitle{Step 4: Data Quality Check}
\begin{block}{Intuition}
No missing values or duplicates here, but the churn split (about 86/14) is a real finding,
not a non-issue: with churners this rare, a model that just always predicts ``stays'' would
already be 85\% ``accurate'' while being completely useless -- a pattern to watch for once
we get to evaluation metrics.
\end{block}
\begin{lstlisting}
df.isnull().sum().sum()
0

df.duplicated().sum()
0

df['Churn'].value_counts(normalize=True)
False    0.855
True     0.145
Name: Churn, dtype: float64
\end{lstlisting}
\end{frame}

%%%%%%%%%%%%%%%%%%%%%%%%%%%%%%%%%%%%%%%%%%%%%%%%%%%%%%%%%%
\begin{frame}[fragile]\frametitle{Step 5: Visualize -- Distributions}
Plotting a numeric column shows shape that summary statistics alone can hide (skew,
multiple peaks, a long tail of outliers).
\begin{center}
\includegraphics[width=0.75\linewidth,keepaspectratio]{churn20}
\end{center}
\end{frame}

%%%%%%%%%%%%%%%%%%%%%%%%%%%%%%%%%%%%%%%%%%%%%%%%%%%%%%%%%%
\begin{frame}[fragile]\frametitle{Step 5: Visualize -- Grouped by Target}
Splitting a distribution by the target column is where visualization turns into real signal
-- do churners and non-churners actually look different?
\begin{center}
\includegraphics[width=0.75\linewidth,keepaspectratio]{churn21}
\end{center}
\end{frame}

%%%%%%%%%%%%%%%%%%%%%%%%%%%%%%%%%%%%%%%%%%%%%%%%%%%%%%%%%%
\begin{frame}[fragile]\frametitle{Step 6: Feature Engineering -- New Columns}
Raw columns rarely capture the most useful signal directly -- combine them into something
more meaningful.
\begin{itemize}
\item \texttt{Total minutes} collapses three columns into one overall usage signal.
\item \texttt{High service calls} turns a raw count into a flag -- customers who call
  support a lot are often frustrated customers.
\end{itemize}
\begin{lstlisting}
df['Total minutes'] = (df['Total day minutes'] + df['Total eve minutes']
                        + df['Total night minutes'])

df['High service calls'] = df['Customer service calls'] >= 4
\end{lstlisting}
\end{frame}

%%%%%%%%%%%%%%%%%%%%%%%%%%%%%%%%%%%%%%%%%%%%%%%%%%%%%%%%%%
\begin{frame}[fragile]\frametitle{Step 6: Feature Engineering -- Encoding}
\texttt{International plan} and \texttt{Voice mail plan} are Yes/No text -- most models need
numbers.
\begin{lstlisting}
df['International plan'] = (df['International plan'] == 'Yes').astype(int)
df['Voice mail plan'] = (df['Voice mail plan'] == 'Yes').astype(int)

df[['International plan', 'Voice mail plan']].head(3)
   International plan  Voice mail plan
0                    0                1
1                    0                1
2                    0                0
\end{lstlisting}
\end{frame}

%%%%%%%%%%%%%%%%%%%%%%%%%%%%%%%%%%%%%%%%%%%%%%%%%%%%%%%%%%
\begin{frame}[fragile]\frametitle{Step 7: What Actually Matters? -- Correlation}
\begin{block}{Intuition}
None of these correlations is huge on its own -- that's normal and expected. It's a first
pass, not a final answer: it tells you where to look closer (service calls, day-time usage,
having an international plan) before committing to a full model in the next session.
\end{block}
\begin{lstlisting}
df.corr(numeric_only=True)['Churn'].sort_values(ascending=False)
Customer service calls    0.21
High service calls        0.19
Total day minutes         0.21
International plan        0.26
Voice mail plan          -0.10
\end{lstlisting}
\end{frame}

%%%%%%%%%%%%%%%%%%%%%%%%%%%%%%%%%%%%%%%%%%%%%%%%%%%%%%%%%%
\begin{frame}[fragile]\frametitle{Step 7: What Actually Matters? -- Group Comparison}
Correlation only captures straight-line relationships -- groupby often reveals more. Here,
churners call support noticeably more on average -- exactly the kind of pattern a model can
learn from.
\begin{lstlisting}
df.groupby('Churn')['Customer service calls'].mean()
Churn
False    1.45
True     2.23
Name: Customer service calls, dtype: float64
\end{lstlisting}
\end{frame}

%%%%%%%%%%%%%%%%%%%%%%%%%%%%%%%%%%%%%%%%%%%%%%%%%%%%%%%%%%
\begin{frame}[fragile]\frametitle{Recap: The Pattern}
\begin{itemize}
\item \textbf{Load} -- get the data into a DataFrame.
\item \textbf{Assess} -- check shape, columns, types.
\item \textbf{Describe} -- summary statistics, missing values, duplicates, class balance.
\item \textbf{Visualize} -- distributions, and distributions split by the target.
\item \textbf{Engineer} -- new columns, encode categoricals.
\item \textbf{Check what matters} -- correlation and groupby comparisons.
\end{itemize}
This exact loop is what the rest of this course's ML sessions assume you've already done
before a single model gets trained.
\end{frame}

%%%%%%%%%%%%%%%%%%%%%%%%%%%%%%%%%%%%%%%%%%%%%%%%%%%%%%%%%%
\begin{frame}[fragile]\frametitle{Quick Check: EDA \& Data Preparation}
\begin{block}{Think About It}
The churn split was about 86\% stay / 14\% churn. Why does that number matter more for
choosing an \emph{evaluation metric} later than it does for the EDA and preparation steps
we just did?
\end{block}
\pause
\begin{block}{Answer}
Nothing in loading, cleaning, or engineering features depends on the class balance -- those
steps work the same either way. But once you train a model, plain accuracy becomes
misleading with an 86/14 split, since always predicting ``stays'' already scores 86\%. EDA's
job here is to \emph{surface} the imbalance early; deciding what to do about it (a different
metric, resampling, etc.) is a modeling-stage decision, covered later in this course.
\end{block}
\end{frame}
